% Replacement for the supplied Results, Discussion, and Limitations sections.
% Uses the same frozen 145-question run; post-hoc scores do not regenerate answers.
% Required packages: graphicx, float, amsmath.

\section{Results and Analysis}
\fontsize{11}{13}\selectfont

\subsection{Evaluation Scope}
The four configurations were evaluated on the same 145 held-out FinanceBench
questions. P0 generated answers without retrieval; retrieval metrics therefore
apply only to P1--P3. The abstention analysis and the additional correctness,
reference-coverage and BERTScore evaluations were performed on the saved answers
from this common run. They did not generate new answers or change the recorded
retrieval contexts and response times.

The evaluation distinguishes retrieval of annotated pages, support for generated
factual units, whole-answer correctness, coverage of reference answer units,
textual similarity, and exact abstention. These measures have different
denominators and meanings. A response can be supported by its retrieved context
while failing to answer the question correctly; a response can also avoid the
prescribed abstention phrase while providing no substantive answer.

\subsection{Main Results}
\subsubsection{Retrieval of Annotated Evidence}
\begin{table}[H]
\centering
\caption{Retrieval of annotated FinanceBench evidence pages by the three RAG
configurations. Scores are question-level means over all 145 evaluation
questions, based on the final ranked chunks supplied to the answer model.
``Empty'' counts questions with no retrieved chunks; these questions receive
zero retrieval scores.}
\label{tab:main-retrieval-results}
\resizebox{\columnwidth}{!}{%
\begin{tabular}{|c|c|c|c|c|c|c|}
\hline
\textbf{ID} & \textbf{Hit@1} & \textbf{Hit@3} & \textbf{Hit@5}
& \textbf{Recall@5} & \textbf{MRR@5} & \textbf{Empty} \\
\hline
P1 & 0.110 & 0.303 & 0.400 & 0.362 & 0.221 & 2 \\
P2 & 0.097 & 0.234 & 0.317 & 0.290 & 0.174 & 2 \\
P3 & 0.214 & 0.352 & 0.400 & 0.369 & 0.284 & 2 \\
\hline
\end{tabular}%
}
\end{table}
P3 placed an annotated evidence page first most often, with Hit@1 of 0.214 and
MRR@5 of 0.284. P1 and P3 nevertheless had identical Hit@5 of 0.400: each
retrieved at least one annotated page on 58 questions. Their Recall@5 means
differed by only 0.007. P2 scored below P1 on every reported retrieval metric.
These results distinguish earlier placement of annotated evidence from finding
that evidence anywhere in the final five chunks. They do not establish whether
the selected chunk contains the precise answer-bearing passage or table cells.

\subsubsection{Whole-Answer Correctness and Reference Coverage}
A post-hoc evaluation used \texttt{gpt-4o-mini} to generate 145 reference-unit rubrics from the questions,
reference answers and annotated FinanceBench evidence, without seeing any
pipeline answers. One rubric per question was shared across P0--P3. The 179
exact abstentions were assigned deterministic failure-to-answer outcomes and
zero correctness and coverage scores. The remaining 401 responses were
submitted to \texttt{gpt-4o-mini} for comparison with the question, reference
answer, annotated evidence and shared rubric. The judge did not receive the
pipeline identity or the retrieved context.
Whole-answer correctness was binary: a valid judgment scored one only when all
required rubric units were answered and no substantive contradiction was
identified. Other valid outcomes scored zero. Partial answers were retained as
a separate outcome category, rather than counted as fully correct.

The additional completeness measure was the fraction of generated reference
units judged correctly answered, macro-averaged across questions. It is reported
as a \emph{reference-unit coverage proxy}. The intended rubric contained only
essential requested information, but inspection found that some rubrics also
required optional supporting quantities. This measure therefore cannot be
assumed to capture only the minimum content needed to answer the question.
An answer can receive some coverage credit while remaining incorrect because
of a wrong requested result or contradictory reasoning.

\begin{table}[H]
\centering
\caption{Post-hoc model-judged whole-answer correctness and reference-unit
coverage against the question, reference answer and annotated gold evidence,
with BERTScore similarity to the reference answer. Correctness and coverage
are means over the valid semantic judgments shown in $n$; exact abstentions
are valid zero scores. Judge errors have missing semantic scores. BERTScore
F1 uses all 145 answers per configuration, including abstentions.}
\label{tab:main-semantic-results}
\resizebox{\columnwidth}{!}{%
\begin{tabular}{|c|c|c|c|c|c|}
\hline
\textbf{ID} & \textbf{Valid $n$} & \textbf{Judge errors}
& \textbf{Correctness} & \textbf{Coverage proxy} & \textbf{BERTScore F1} \\
\hline
P0 & 143 & 2 & 0.098 & 0.135 & 0.475 \\
P1 & 142 & 3 & 0.268 & 0.299 & 0.467 \\
P2 & 141 & 4 & 0.220 & 0.256 & 0.466 \\
P3 & 143 & 2 & 0.266 & 0.300 & 0.476 \\
\hline
\end{tabular}%
}
\end{table}

Of the 401 responses submitted to the judge, 390 produced valid judgments.
Together with the 179 deterministic abstention scores, these formed 569 valid
semantic results across the 580 responses. The remaining 11 judgments failed
consistency validation after three attempts and were retained as unresolved
judge errors, with no imputed semantic score.
Table~\ref{tab:answer-outcomes} gives the full outcome counts. P1 and P3 each
had 38 answers judged fully correct, compared with 31 for P2 and 14 for P0.
P1's slightly higher valid-judgment correctness mean than P3's arises from the
different denominators, 142 and 143, rather than an additional correct answer.

\begin{table}[H]
\centering
\caption{Whole-answer outcome counts from the post-hoc reference-based
evaluation. The five outcome columns partition all 145 responses per
configuration. Failure to answer includes prescribed exact abstentions and
broader responses judged not to supply a substantive answer. An unresolved
judge error is a scoring failure, not an incorrect-answer judgment.}
\label{tab:answer-outcomes}
\resizebox{\columnwidth}{!}{%
\begin{tabular}{|c|c|c|c|c|c|}
\hline
\textbf{ID} & \textbf{Correct} & \textbf{Partially correct}
& \textbf{Incorrect} & \textbf{Failure to answer} & \textbf{Judge error} \\
\hline
P0 & 14 & 8 & 114 & 7 & 2 \\
P1 & 38 & 6 & 40 & 58 & 3 \\
P2 & 31 & 7 & 39 & 64 & 4 \\
P3 & 38 & 5 & 42 & 58 & 2 \\
\hline
\end{tabular}%
}
\end{table}

The uncertainty introduced specifically by unresolved judgments can be made
explicit over the full 145-question denominator. If $c$ is the number judged
correct and $u$ is the number unresolved, the lower and upper correctness
bounds are $c/145$ and $(c+u)/145$. These bounds are 0.097--0.110 for P0,
0.262--0.283 for P1, 0.214--0.241 for P2, and 0.262--0.276 for P3. They are
not confidence intervals and do not account for undetected judge mistakes.

Because the RAG configurations abstained on different questions, the
conditional results are also reported in
Table~\ref{tab:conditional-answer-results}. These exclude exact abstentions;
semantic means additionally exclude unresolved judge errors. A longer
procedural response or noncanonical refusal remains in this non-exact-
abstention population, including when it is judged a failure to answer.

\begin{table}[H]
\centering
\caption{Answer evaluation conditional on not producing a prescribed exact
abstention. ``Responses'' counts the entire non-exact-abstention subset;
``Valid semantic $n$'' excludes unresolved judge errors. Correctness and
reference-unit coverage use the latter denominator. BERTScore uses all
responses in this subset. The selected question sets differ across pipelines.}
\label{tab:conditional-answer-results}
\resizebox{\columnwidth}{!}{%
\begin{tabular}{|c|c|c|c|c|c|}
\hline
\textbf{ID} & \textbf{Responses} & \textbf{Valid semantic $n$}
& \textbf{Correctness} & \textbf{Coverage proxy} & \textbf{BERTScore F1} \\
\hline
P0 & 145 & 143 & 0.098 & 0.135 & 0.475 \\
P1 & 87 & 84 & 0.452 & 0.506 & 0.557 \\
P2 & 82 & 78 & 0.397 & 0.463 & 0.568 \\
P3 & 87 & 85 & 0.447 & 0.505 & 0.572 \\
\hline
\end{tabular}%
}
\end{table}

Conditional correctness was 0.452 for P1 and 0.447 for P3, compared with
0.397 for P2. These means describe selected responses, rather than performance
over all 145 questions. They should be read alongside abstention and scoring
coverage, rather than used alone to rank the configurations.

BERTScore used \texttt{bert-base-uncased}, layer 9, without IDF weighting or
baseline rescaling. All 580 answer--reference pairs were scored without
truncation. Its F1 measures contextual token similarity, not the proportion of
fully correct answers. P3 had the highest overall BERTScore F1, 0.476, followed
closely by P0 at 0.475, despite their much larger difference in model-judged
correctness. Thus, similarity and correctness yield different descriptions
of these responses. BERTScore values also depend on the encoder and scoring
configuration and should not be directly compared with results using another
variant.

These semantic scores remain exploratory model judgments. A qualitative audit
found that P0's statement of ``over 60 years'' of dividend increases was
rejected against evidence of 65 years, although the statements are compatible
(\texttt{financebench\_id\_01858}). Another answer correctly stated that
Adobe's free cash-flow conversion improved but was penalised for omitting
supporting percentages (\texttt{financebench\_id\_00591}). These observations
show that schema consistency does not establish semantic validity. The audit
was a convenience inspection by the code assistant, not independent human
validation or an estimate of the judge's error rate. No answer scores were
manually changed using these observations.

\subsubsection{Lexical and Numeric Diagnostics}
\begin{table}[H]
\centering
\caption{Lexical and numeric comparison with the FinanceBench reference
answers. Exact match (EM) and token F1 use all 145 questions. Numeric metrics
use the 112 questions with extractable reference numbers and the corrected
numeric extractor. ``Numeric accuracy'' is the question-level proportion of
reference numeric entities matched, rather than a whole-answer correctness rate.}
\label{tab:main-answer-results}
\resizebox{\columnwidth}{!}{%
\begin{tabular}{|c|c|c|c|c|c|}
\hline
\textbf{ID} & \textbf{EM} & \textbf{Token F1}
& \textbf{Numeric accuracy} & \textbf{Numeric precision}
& \textbf{Numeric F1} \\
\hline
P0 & 0.007 & 0.172 & 0.056 & 0.051 & 0.048 \\
P1 & 0.014 & 0.159 & 0.182 & 0.123 & 0.130 \\
P2 & 0.007 & 0.166 & 0.190 & 0.130 & 0.137 \\
P3 & 0.007 & 0.170 & 0.176 & 0.114 & 0.127 \\
\hline
\end{tabular}%
}
\end{table}
The three RAG pipelines had higher numeric-entity recall than P0, with P2
highest at 0.190. Exact match remained low: P1 matched two reference answers
exactly, while each other configuration matched one. These original
diagnostics were preserved when the semantic and BERTScore columns were added.
P2's highest numeric score, alongside its lower model-judged correctness than
P1 and P3, illustrates why matching reference numbers does not establish
correct interpretation, scope or reasoning for a complete answer.

\subsubsection{Factual-Unit Support and Hallucination Proxies}
\begin{table}[H]
\centering
\caption{Model-judged support for generated factual units and complementary
hallucination proxies. Context refers to the final retrieved chunks; gold
refers to annotated FinanceBench evidence. Parentheses give the number of
scored answers. Unsupported units are not independently verified false claims.}
\label{tab:main-support-results}
\resizebox{\columnwidth}{!}{%
\begin{tabular}{|c|c|c|c|c|c|}
\hline
\textbf{ID} & \textbf{Exact abstentions} & \textbf{Context support}
& \shortstack{\textbf{Context halluc.}\\\textbf{proxy}}
& \textbf{Gold support} & \shortstack{\textbf{Gold halluc.}\\\textbf{proxy}} \\
\hline
P0 & 0  & NA          & NA    & 0.583 (112) & 0.417 \\
P1 & 58 & 0.850 (85)  & 0.150 & 0.636 (83)  & 0.364 \\
P2 & 63 & 0.820 (80)  & 0.180 & 0.621 (81)  & 0.379 \\
P3 & 58 & 0.832 (86)  & 0.168 & 0.652 (85)  & 0.348 \\
\hline
\end{tabular}%
}
\end{table}
For each answer with at least one judge-scored factual unit, support was the
proportion of those units judged supported by the selected evidence source.
The corresponding hallucination proxy was $1-\text{support}$, the proportion
judged unsupported. Values are question-level macro means over the same
applicable answers for each support--proxy pair. Consequently, each pair sums
to one by definition: P1's context support of 0.850 and proxy of 0.150 are two
views of the same judgments, rather than independent findings. Context and
gold support use different evidence sources and should not be added together.

P1 had the highest conditional context-support mean, while P3 had the highest
conditional gold-support mean. These results concern the generated factual
units that were scored, rather than the number of questions answered
correctly. For example, P1's context mean covers 85 answers, fewer than both
its 145 evaluation questions and its 87 non-abstentions. Its post-hoc
whole-answer correctness uses 142 valid semantic results, including its 58
exact abstentions as zeros. These denominators describe different populations
and must not be interchanged. P0 has no context score because it received no
retrieved chunks.

\subsubsection{Latency}
\begin{table}[H]
\centering
\caption{Per-question response latency in milliseconds ($n=145$ per
configuration). Total time includes retrieval, reranking and answer
generation where applicable; median and standard deviation (SD) describe
total time. NA marks stages a configuration does not use. These timings
exclude startup and the subsequent evaluation calls.}
\label{tab:main-latency-results}
\resizebox{\columnwidth}{!}{%
\begin{tabular}{|c|r|r|r|r|r|r|}
\hline
\textbf{ID} & \textbf{Total mean} & \textbf{Median}
& \textbf{Total SD} & \textbf{Retrieval mean}
& \textbf{Rerank mean} & \textbf{Generation mean} \\
\hline
P0 & 2,074 & 1,916 & 1,006 & NA  & NA    & 2,074 \\
P1 & 1,933 & 1,795 & 853   & 86  & NA    & 1,847 \\
P2 & 2,246 & 2,176 & 777   & 461 & NA    & 1,785 \\
P3 & 5,633 & 5,619 & 903   & 466 & 3,366 & 1,800 \\
\hline
\end{tabular}%
}
\end{table}
P1 had the lowest observed mean total latency at 1,933\,ms, followed by P0 at
2,074\,ms and P2 at 2,246\,ms. P3 averaged 5,633\,ms, including 3,366\,ms
recorded as reranking. The post-hoc correctness and similarity evaluation
does not alter these response-time measurements.

\subsection{Abstention Behaviour}
\subsubsection{Frequency, Shared Questions and Overlap}
Across P1--P3, there were 179 exact abstention decisions among 435 RAG
responses (41.1\%). They affected 77 distinct questions (53.1\% of 145).
These totals count different units: 179 counts pipeline responses, whereas
77 counts questions with at least one abstaining pipeline. P1 and P3 each
abstained on 58/145 questions (40.0\%), leaving 87 responses without an exact
abstention. P2 abstained on 63/145 (43.4\%), leaving 82. P0 recorded no exact
abstentions. Including P0 gives 179/580 responses (30.9\%), although the
RAG-only rate is more directly relevant to refusal under the retrieved-context
restriction.

The frozen exact-abstention detector checks whether the entire normalised
response matches one of five prescribed refusal phrases. All 179 detected
responses used ``Insufficient evidence in the retrieved context.'' Longer
refusals, procedural explanations and incomplete answers can remain outside
this count. Table~\ref{tab:answer-outcomes} demonstrates the distinction:
P0 had seven model-judged failures to answer despite zero exact abstentions,
and P2 had 64 failures to answer compared with 63 exact abstentions.
These broader labels are themselves subject to judge error.

\begin{table}[H]
\centering
\caption{Overlap of exact-abstention decisions across the three RAG
configurations on the same 145 questions. ``Only'' excludes abstention by
the remaining configuration or configurations.}
\label{tab:abstention-overlap}
\begin{tabular}{|l|r|r|}
\hline
\textbf{Abstaining configurations} & \textbf{Questions} & \textbf{Percentage} \\
\hline
None                 & 68  & 46.9\% \\
P1 only              & 4   & 2.8\% \\
P2 only              & 7   & 4.8\% \\
P3 only              & 7   & 4.8\% \\
P1 and P2 only       & 8   & 5.5\% \\
P1 and P3 only       & 3   & 2.1\% \\
P2 and P3 only       & 5   & 3.4\% \\
P1, P2 and P3        & 43  & 29.7\% \\
\hline
Total                & 145 & 100.0\% \\
\hline
\end{tabular}
\end{table}

All three RAG configurations abstained on 43 questions, representing 55.8\%
of the 77 questions with any abstention. These shared cases accounted for
74.1\% of P1's abstentions, 68.3\% of P2's and 74.1\% of P3's. No pipeline
exact-abstained on 68 questions; the remaining 34 had mixed decisions.
Exactly one pipeline abstained on 18 questions and exactly two on 16.

P1 and P2 had the most similar abstention sets: 51 shared abstentions and
Jaccard overlap of 72.9\%. P1--P3 shared 46 (65.7\%), and P2--P3 shared
48 (65.8\%). Here, Jaccard overlap divides shared abstentions by questions
where either pipeline abstained. Equal totals did not mean identical behaviour:
P1 and P3 changed status on 24 questions, with 12 transitions in each
direction. P3 therefore redistributed exact abstentions relative to P1
without reducing their total. Relative to P2, P3 changed 15 abstentions into
non-abstentions but introduced ten elsewhere, giving a net reduction of five.

\subsubsection{Evidence Availability and Selection Stages}
Among abstentions, a final annotated gold-page hit was absent on 49/58 P1
responses, 57/63 P2 responses and 48/58 P3 responses. Together, these account
for 154/179 abstention decisions (86.0\%).
Table~\ref{tab:abstention-stages} partitions the exact abstentions by the
observed retrieval stage at which annotated-page coverage was absent or lost.

\begin{table}[H]
\centering
\caption{Observed annotated-page retrieval conditions among exact
abstentions. These mutually exclusive stage conditions partition the 179
abstaining RAG responses; they are not validated internal reasons for refusal.
Candidate-pool absence includes the two empty retrievals per configuration.}
\label{tab:abstention-stages}
\resizebox{\columnwidth}{!}{%
\begin{tabular}{|l|r|r|r|r|}
\hline
\textbf{Observed condition} & \textbf{P1} & \textbf{P2} & \textbf{P3} & \textbf{Total} \\
\hline
Gold page absent from observed candidate pool & 49 & 40 & 29 & 118 \\
Gold page present in pool but lost at RRF cutoff & 0 & 17 & 7 & 24 \\
Gold page retained before reranking but absent from final five & 0 & 0 & 12 & 12 \\
At least one gold page present in final context & 9 & 6 & 10 & 25 \\
\hline
Total exact abstentions & 58 & 63 & 58 & 179 \\
\hline
\end{tabular}%
}
\end{table}

The observed pools differ by architecture. P1's dense top five also formed its
final context. P2 combined dense and sparse top-ten results before selecting
the RRF top five. P3 combined dense and sparse top-twenty results, retained
the RRF top twenty and selected the reranked top five. Consequently, these
diagnostics describe complete configurations with different search budgets.
They do not isolate the effect of reranking alone. Each pipeline returned
an empty context on the same two questions; most refusals therefore occurred
with some retrieved material available. Empty retrieval is included in
Table~\ref{tab:abstention-stages} and must not be added to its totals.

Abstention was more frequent without a final gold-page hit. Conditional on
having no hit, the rates were 49/87 (56.3\%) for P1, 57/99 (57.6\%) for P2
and 48/87 (55.2\%) for P3. With a hit, the rates were 9/58 (15.5\%),
6/46 (13.0\%) and 10/58 (17.2\%), respectively. These denominators contain
all questions with the relevant retrieval condition, rather than only
abstentions. They show an association between annotated evidence availability
and refusal, without establishing an individual model's internal reason.

Among the 43 jointly abstained questions, 35 had no final gold-page hit in
any pipeline, three had hits in all three and five had mixed retrieval
outcomes. Conversely, 22 of the 68 questions without any exact abstention
had no final gold-page hit in any pipeline. Neither condition determines
answer correctness: unannotated passages may provide valid alternative
evidence, while a chunk from an annotated page may omit the decisive passage.

\subsubsection{Abstention Despite Retrieval of an Annotated Page}
There were 25 exact abstention decisions despite at least one annotated
gold page being present in the final context, affecting 16 distinct questions.
These comprised nine P1 decisions, six P2 decisions and ten P3 decisions:
15.5\%, 9.5\% and 17.2\% of each pipeline's abstentions, respectively,
or 14.0\% of all 179 abstentions. These percentages differ from the
conditional rates above because the denominators here are abstentions.

Ten of the 25 decisions, across six questions, retrieved only one of two
annotated evidence pages. The remaining 15 decisions concerned questions
with a single annotated page, which was retrieved. This verifies page
coverage, rather than sufficient answer-bearing passage coverage. Inspection
of selected saved contexts illustrates several distinct situations:

\begin{itemize}
\item \textbf{Incomplete evidence across pages.} For Adobe's operating
cash-flow ratio (\texttt{financebench\_id\_03856}), P3 retrieved the cash-flow
page but not the annotated balance-sheet page containing total current
liabilities. Retrieving one correct page therefore left a required input
missing.
\item \textbf{A correct page with an incomplete table fragment.} For Foot
Locker's board-election question (\texttt{financebench\_id\_00822}), P1's
selected chunk began partway through the voting table, omitting the headings
and Richard A. Johnson's row. P2 retrieved both relevant chunks and identified
Richard A. Johnson. The page hit alone overstated the passage coverage
available to P1.
\item \textbf{Refusal despite explicit answer-bearing evidence.} For AMD's
segment-growth question (\texttt{financebench\_id\_00563}), P3's final context
explicitly stated that Data Center revenue increased by 64\%, supporting the
answer produced by P1 and P2. P3 nevertheless exact-abstained. Its other four
chunks concerned FY2015. Distraction by older material is plausible, but the
saved refusal does not establish that cause.
\item \textbf{A negative inference from a listing.} For Ulta's registered
debt securities (\texttt{financebench\_id\_00746}), all three pipelines
abstained despite retrieving the annotated cover-page listing. The annotated
answer, ``There are none,'' required interpreting the listing, which showed
common stock, rather than copying an explicit affirmative debt entry.
\end{itemize}

These cases distinguish \emph{abstention despite an annotated-page hit} from
\emph{abstention despite sufficient answer-bearing evidence}. The first has
the verified count of 25 decisions across 16 questions; AMD supplies a clear
example of the second. The saved data do not provide validated, mutually
exclusive reason counts for all 25 decisions, so the examples should not be
presented as an exhaustive causal classification.

An evidence-selection example also helps explain differences across pipelines.
For Best Buy's FY2019 inventories
(\texttt{financebench\_id\_04417}), P1 and P2 abstained, while P3 answered
``\$5,409 million,'' matching the gold value. P1's observed candidate pool
lacked the annotated page; P2 found an annotated-page candidate but lost it
at the RRF cutoff; P3 retained a final chunk explicitly listing merchandise
inventories of 5,409 million. This is an observed correspondence between
evidence surviving selection and the response decision, without establishing
what would happen under a controlled change to one component.

\subsubsection{Question Types and Broader Failures to Answer}
\begin{table}[H]
\centering
\caption{Exact abstentions by native FinanceBench question type. Each
percentage uses the number of questions in its row. These are descriptive
subgroup results, not additional hypothesis tests.}
\label{tab:abstention-question-types}
\begin{tabular}{|l|r|r|r|r|}
\hline
\textbf{Question type} & \textbf{$n$} & \textbf{P1} & \textbf{P2} & \textbf{P3} \\
\hline
Domain-relevant & 48 & 19 (39.6\%) & 19 (39.6\%) & 17 (35.4\%) \\
Metrics-generated & 48 & 24 (50.0\%) & 28 (58.3\%) & 26 (54.2\%) \\
Novel-generated & 49 & 15 (30.6\%) & 16 (32.7\%) & 15 (30.6\%) \\
\hline
\end{tabular}
\end{table}

Metrics-generated questions had the highest abstention rate in every
configuration, while novel-generated questions had the lowest. P2's five
additional abstentions relative to P1 comprised four metrics-generated and
one novel-generated question. P3's equal total relative to P1 concealed two
fewer domain-relevant abstentions and two more metrics-generated abstentions.

For questions labelled purely Numerical reasoning, the rates were 22/43
(51.2\%), 23/43 (53.5\%) and 25/43 (58.1\%) for P1--P3. For pure
Information extraction, they were 7/29 (24.1\%), 8/29 (27.6\%) and 5/29
(17.2\%). These differences are consistent with difficulty obtaining and
combining numerical inputs, but the labels do not establish a cause. Mixed
reasoning labels form separate categories, and 49 questions lacked reasoning
annotations.

Exact non-abstention should not be treated as successful completion. For MGM
Resorts' average capex-to-revenue question
(\texttt{financebench\_id\_03849}), P1 and P3 exact-abstained. P2 instead
listed missing inputs and ended with the insufficient-evidence phrase,
without providing the requested 7.9\% answer. Because the whole response did
not exactly match a prescribed phrase, it was outside the exact-abstention
count. The added whole-answer evaluation provides a separate, model-judged
assessment of such broader failures.

Abstention also restricts the interpretation of the support results. Beyond
the exact abstentions, P1 had two context-unscored and four gold-unscored
non-abstentions; P2 had two and one; P3 had one and two. These exclusions are
evidence-source-specific and must not be added as distinct questions. P0 had
33 responses without scored gold factual units despite no exact abstentions.
Consequently, the highest conditional support mean does not establish the
best supported response across the complete question set. The response sets
and the further scored subsets both differ.

\subsection{Paired Comparisons}
Paired comparisons aligned answers by FinanceBench question ID and calculated
each difference as configuration B minus A. Binary outcomes, including Hit@5
and exact match, used two-sided exact McNemar tests. Graded scores and latency
used two-sided Wilcoxon signed-rank tests, excluding zero differences from
ranking. Missing or inapplicable scores were excluded pairwise rather than
imputed as zero. Seeded paired bootstrap resampling with 10,000 replicates
produced 95\% percentile intervals for mean paired differences. Holm
adjustment was applied across the 21 existing contrasts.

These contrasts were specified after inspection of earlier aggregate results,
so the inferential analysis is exploratory rather than blinded or
preregistered. Bootstrap intervals estimate mean differences, whereas
Wilcoxon tests signed ranks; their thresholds need not agree. The new
post-hoc semantic, coverage and BERTScore metrics and the abstention subgroups
were not added to this test family. Their comparisons above are descriptive;
the reported adjusted $p$-values do not apply to them.

\begin{table}[H]
\centering
\caption{Selected existing paired comparisons. Mean differences are B minus
A; positive latency differences indicate slower responses. Pairs are
questions with valid scores for both configurations. Reported $p$-values are
Holm-adjusted across the 21 existing comparisons, excluding the newly added
semantic and BERTScore metrics.}
\label{tab:paired-comparisons}
\resizebox{\columnwidth}{!}{%
\begin{tabular}{|l|l|r|r|c|}
\hline
\textbf{A $\rightarrow$ B} & \textbf{Metric}
& \textbf{Mean difference} & \textbf{Pairs} & \textbf{Holm $p$} \\
\hline
P0 $\rightarrow$ P1 & Gold support & $+0.083$ & 66 & 1.000 \\
P0 $\rightarrow$ P1 & Numeric accuracy & $+0.126$ & 112 & 0.030 \\
P0 $\rightarrow$ P3 & Numeric accuracy & $+0.121$ & 112 & 0.028 \\
P1 $\rightarrow$ P2 & Hit@5 & $-0.083$ & 145 & 0.177 \\
P1 $\rightarrow$ P2 & Recall@5 & $-0.072$ & 145 & 0.184 \\
P2 $\rightarrow$ P3 & Context support & $+0.013$ & 71 & 1.000 \\
P2 $\rightarrow$ P3 & Total latency (ms) & $+3{,}387$ & 145 & $<0.001$ \\
\hline
\end{tabular}%
}
\end{table}

For P2 versus P3, the paired context-support difference was $+0.013$
(95\% bootstrap interval $[-0.065,\,0.091]$; Holm-adjusted $p=1.000$).
The mean latency difference was $+3{,}387$\,ms (95\% bootstrap interval
$[3{,}216,\,3{,}554]$\,ms; adjusted $p<0.001$). Thus, the recorded run
established a substantial latency difference between these complete
configurations without establishing a context-support improvement.

The selected numeric comparisons detected improvements for P1 and P3 over
P0, while the paired gold-support comparison did not establish P1's proposed
gain. P2's lower Hit@5 and Recall@5 than P1 were not significant after
correction. None of these results supplies an inferential test of the new
whole-answer correctness differences.

\section{Discussion}
\fontsize{11}{13}\selectfont

The added answer evaluation extends the comparison beyond matching words and
numbers, but it does not replace independent verification of correctness.
P0 had the lowest descriptive model-judged correctness and reference-unit
coverage. P1 and P3 each had 38 answers judged fully correct, and P2 had 31.
This pattern is consistent with a benefit from retrieval in this run, but the
new metrics were evaluated post hoc with an imperfect automated judge.
The original claim that whole-answer correctness was entirely unassessed
must therefore be revised: it is now assessed by a model, but remains
unvalidated by independent human review. The prediction of P0 being least
accurate has descriptive support under this protocol, rather than a confirmed
general ranking of factual correctness.

P2's highest numeric-entity recall did not translate into the highest semantic
correctness or coverage. Likewise, P3's highest BERTScore did not produce more
answers judged fully correct than P1. These results show why numeric matching,
reference similarity, factual-unit support and whole-answer correctness should
be reported separately. In particular, a hallucination proxy derived from
unsupported factual units is not a verified false-answer rate, and its
complementary support score cannot be interpreted as answer accuracy.

Abstention is a major part of the observed system behaviour. More than two
fifths of the RAG responses were exact refusals, and most abstentions within
each pipeline came from the 43 questions jointly abstained by all three.
Missing annotated evidence was common and associated with higher refusal
rates. Nevertheless, pipeline-dependent decisions affected 34 questions,
and equal P1/P3 abstention totals concealed 24 changed decisions. The
configurations therefore differed in which questions they attempted, as well
as in the support and correctness of the resulting responses. Conditional
quality means alone cannot resolve that trade-off.

The 25 refusals despite annotated-page retrieval further separate finding a
relevant page from obtaining and using sufficient evidence. Some contexts
contained only one page of a multi-page calculation or an incomplete table
fragment; AMD supplied an example where explicit answer-bearing evidence was
available but the generator still refused. Possible remedies include retaining
related table fragments, preserving multi-page inputs and reducing irrelevant
context, but these were not experimentally tested here. The refusal text
does not reveal a model's internal cause, so the observed stage conditions
should not be presented as confirmed explanations for every refusal.

The prediction that P0 would be fastest was not supported descriptively:
P1's mean latency was lower, although the paired P0--P1 difference was not
statistically established. P0 also had the highest descriptive gold-evidence
hallucination proxy, but the paired gold-support tests did not establish a
corresponding inferential claim. P1's proposed gold-support advantage over
P0 and P3's proposed context-support advantage over P2 were not established.
A nonsignificant result is not evidence that the configurations perform equally.

P3's stronger first-rank retrieval came with substantially greater latency
and no additional answers judged fully correct relative to P1. P1 and P3
also matched on Hit@5 and exact-abstention totals. For these frozen
implementations, this supports P1 as a provisional latency-conscious baseline.
It does not establish that P1 is the most factually correct architecture.
P3 searched a wider candidate pool than P2 before reranking, so the P2--P3
comparison cannot isolate the cross-encoder's causal contribution. Similarly,
the native question-type and reasoning results are descriptive and do not
establish the earlier terminology-heavy or numeric-heavy subgroup predictions.

\section{Limitations and Conclusion}
\fontsize{11}{13}\selectfont

Several limitations affect the interpretation of these results. Retrieval was
evaluated against annotated pages rather than answer-bearing passages or
table cells. Annotated-page absence can miss alternative valid evidence,
and annotated-page presence can conceal insufficient passage coverage.
The exact-abstention detector captures a narrow response form; the additional
model-judged failure-to-answer category broadens the analysis but introduces
its own classification errors. The qualitative refusal examples were selected
for explanation and do not provide an exhaustive, validated causal taxonomy.

The support and whole-answer evaluations used \texttt{gpt-4o-mini}, the same
model family used for generation in separate calls. No independent human
assessment of answer correctness was performed. Inspection found compatible
phrasing rejected, optional supporting details treated as required, and
period or scope omissions in generated rubrics. Shared rubrics provide a
common standard across architectures but can also propagate the same rubric
defect to all of them. The 11 unresolved judgments remain missing, and the
reported correctness bounds account only for those unresolved cases, not
undetected errors among accepted judgments. Small differences in correctness,
coverage or factual support should therefore be interpreted cautiously.

Reference-unit coverage is a completeness proxy, while EM, token F1,
numeric-entity scores and BERTScore are complementary diagnostics rather than
independent determinations of complete factual correctness. The reported
BERTScore variant uses BERT-base with layer 9 and no baseline rescaling;
its values are not directly interchangeable with another encoder's scores.
Differences in datasets, evidence, judging rubrics and abstention handling
also prevent direct numerical comparison with published systems using other
evaluation protocols.

The study used one common 145-question run, with no repeated generation or
timing runs, tail-latency evaluation or monetary-cost comparison. Response
timings exclude startup and the later evaluation work. Candidate depths
differed between architectures, preventing a controlled attribution of
observed changes to reranking alone. The existing hypothesis tests are
exploratory, and the added semantic and abstention analyses remain descriptive.

For the frozen implementations, P3 ranked annotated evidence earlier but
incurred substantially higher latency without an established context-support
gain or a higher count of answers judged fully correct than P1. P2 did not
improve the observed retrieval means over P1. The abstention analysis reveals
both shared evidence-access difficulty and configuration-dependent response
decisions, including occasional refusal despite explicit answer-bearing
context. P1 is consequently a reasonable provisional baseline when latency
matters, rather than a demonstrated accuracy optimum. A further study should
independently validate correctness and completeness, assess passage-level
evidence sufficiency, and compare reranking at equal candidate depths while
reporting answer coverage alongside conditional quality.
